\documentclass[12pt,letterpaper]{article}

\usepackage[T1]{fontenc}
\usepackage[utf8]{inputenc}
\usepackage{newtxtext,newtxmath}
\usepackage[margin=0.9in,top=0.75in,bottom=0.75in]{geometry}
\usepackage{microtype}
\usepackage{enumitem}
\usepackage{titlesec}
\usepackage[dvipsnames]{xcolor}
\usepackage{xurl}
\usepackage[hidelinks]{hyperref}

\definecolor{ink2}{HTML}{52514E}
\definecolor{rule}{HTML}{C3C2B7}
\pagestyle{empty}
\setlength{\parindent}{0pt}
\setlength{\parskip}{4pt}
\titleformat{\section}{\normalsize\bfseries}{}{0pt}{}
\titlespacing*{\section}{0pt}{7pt}{2pt}
\setlist{itemsep=1pt,topsep=2pt,parsep=0pt,leftmargin=1.6em}

\begin{document}

{\small\color{ink2}COMP 6934 Data Visualization \hfill Project Proposal \hfill Group 16}\\[-4pt]
{\color{rule}\rule{\linewidth}{0.5pt}}

\begin{center}
{\Large\bfseries Less Fish, More Money?}\\[3pt]
{\large Three Decades of Change in Canada's Commercial Sea Fisheries, 1990--2024}\\[6pt]
{\small Utsav Karki (202683489) \quad\textperiodcentered\quad Lata Airee (202680103) \quad\textperiodcentered\quad Abishek Paudel (202684964)}
\end{center}

\section{Dataset}
We will use \emph{Commercial Seafisheries Landings, 1990--2024}, published by Fisheries and
Oceans Canada (DFO) from its Zonal Interchange File database
(\url{https://www.dfo-mpo.gc.ca/stats/commercial/sea-maritimes-eng.htm}). The table has about
10{,}600 records. Each record gives, for one year, one region and one species, the
\textbf{landed quantity} (kg, live weight) and the \textbf{landed value} (CAD paid to harvesters at
the dock). It covers 35 years, the five Atlantic provinces (NL, NS, NB, PE, QC), British
Columbia, and Atlantic and national totals. It includes 33 species, such as cod, lobster, snow
crab, herring and salmon, grouped into four groups: groundfish, pelagics, shellfish and other.
Some provincial species values are suppressed (``x'') for confidentiality, but all group and
national totals are complete. To compare dollars across three decades fairly, we will adjust
values for inflation using the Statistics Canada Consumer Price Index (Table 18-10-0005-01).

\section{Why this dataset}
In 1992 Canada closed the northern cod fishery off Newfoundland and Labrador after the stock
collapsed, one of the best-known fisheries failures in the world. For us, studying at Memorial
University, this history is close to home, because the fishery shaped the economy and
communities of this province. The moratorium is usually remembered only as a loss. This dataset
covers the entire period after the collapse, so we can see what the industry became: which
species replaced cod, whether the fishery recovered in economic terms, and which provinces won
or lost. The questions matter to fisheries managers, coastal communities and anyone interested
in how an economy adapts after a natural resource collapses. Because the data has several
dimensions (time, geography, species, quantity and value), it is well suited to a range of
visualization types.

\section{Questions we intend to investigate}
\begin{enumerate}[label=\textbf{Q\arabic*.}, leftmargin=2.3em]
  \item \textbf{Growth or decline?} Has Canada's sea fishery grown or shrunk since 1990, and does
        the answer change when it is measured in tonnes rather than in inflation-adjusted dollars?
  \item \textbf{What replaced groundfish?} Which species groups and species filled the gap left
        by the collapse of cod and other groundfish, and was the collapse limited to cod?
  \item \textbf{Quantity or price?} Is any growth in value driven by catching more fish, or by
        each kilogram becoming more valuable?
  \item \textbf{Who, and at what risk?} How were these changes distributed across provinces,
        and has the fishery become dependent on a small number of species?
\end{enumerate}

\section{Planned approach}
We will clean and analyse the data in Python (pandas, seaborn, matplotlib and scipy in a Jupyter
notebook). Our plots will include line charts of volume against value indexed to 1990, stacked
area charts of species-group composition, a species-by-year heatmap, a 1990-versus-2024 dumbbell
chart, price-against-quantity scatter plots, provincial stacked bar charts, and box plots of
year-to-year volatility. Each plot will support a specific answer to one of the questions above.

\end{document}
